\chapter{State of the Art}
\label{ch:pfe-sota}

\section{Introduction and Review Scope}

Recovering an SBOM from a compiled Android APK intersects several research
areas: third-party library identification, exact-version detection,
obfuscation-resilient similarity, software composition analysis, and
vulnerability mapping. This chapter reviews the approaches most directly
related to those tasks and identifies the gap addressed by the proposed
system.

The review corpus was assembled from the project paper collection and then
audited against accessible primary PDFs, publisher records, arXiv records, and
public artifacts. The comparison focuses on:

\begin{itemize}
  \item library-level and version-level identification;
  \item the representation used to compare APK and reference code;
  \item resistance to renaming, shrinking, repackaging, and optimization;
  \item partial-library and multidex behavior;
  \item dataset realism and metric interpretation;
  \item runtime, availability, and reproducibility;
  \item standard SBOM output and vulnerability mapping.
\end{itemize}

Reported percentages are not treated as a universal ranking. The reviewed
papers use different datasets, obfuscators, reference libraries, thresholds,
and definitions of a correct match. Original-paper results are therefore
presented as evidence about a method under its own evaluation conditions,
while the later common-benchmark results are used to discuss realistic
generalization.

\section{Evolution of Android Third-Party Library Detection}

Android third-party library detection evolved from fast package and frequency
signatures toward structural correspondence, optimization-aware matching, and
learned representations. Figure~\ref{fig:pfe-sota-evolution} summarizes this
progression. The generations overlap: later tools often retain efficient
signatures as a coarse stage and add deeper evidence only for shortlisted
candidates.

\begin{figure}[H]
\centering
\resizebox{\textwidth}{!}{%
\begin{tikzpicture}[
  era/.style={
    rectangle,
    rounded corners,
    draw=schoolblue,
    thick,
    fill=blue!5,
    align=center,
    text width=42mm,
    minimum height=28mm
  },
  arrow/.style={-{Latex[length=2.5mm]},thick,draw=schoolblue},
  note/.style={font=\scriptsize,align=center,text width=39mm}
]
\node[era] (g1) at (0,0) {
  \textbf{Package and reference signatures}\\
  LibScout, LibRadar, LibD, ORLIS\\[1mm]
  \scriptsize Packages, API frequencies, hashes, CFGs
};
\node[era] (g2) at (5.2,0) {
  \textbf{Dependency and class structure}\\
  LibPecker, LibID, ATVHunter\\[1mm]
  \scriptsize Class dependencies, candidate correspondence, fine opcode checks
};
\node[era] (g3) at (10.4,0) {
  \textbf{Optimization-aware matching}\\
  LibScan, LibHunter, SAD\\[1mm]
  \scriptsize Call chains, optimization simulation, structural version evidence
};
\node[era] (g4) at (15.6,0) {
  \textbf{Learned and SBOM-oriented systems}\\
  LibAttention, UniBOM, this work\\[1mm]
  \scriptsize Learned representations, standard SBOMs, evidence governance
};

\draw[arrow] (g1) -- (g2);
\draw[arrow] (g2) -- (g3);
\draw[arrow] (g3) -- (g4);
\end{tikzpicture}%
}
\caption{Evolution from fragile package signatures toward structural,
optimization-aware, learned, and SBOM-oriented analysis.}
\label{fig:pfe-sota-evolution}
\end{figure}

\subsection{Package, Frequency, and Reference Profiles}

Early approaches sought compact signatures that could be generated and
compared at scale. LibRadar, for example, used Android API frequencies and
package clustering to discover reused code quickly. Its design favored speed,
but clustering did not provide reliable exact-version labels and remained
dependent on package organization. Later systematic evaluations found that
package-structure-based candidate construction is error-prone and that
multidex, partial imports, and library dependencies create additional blind
spots \cite{zhan2020assessment,zhan2021systematic}.

LibScout established a reference-profile model for identifying known
libraries and versions. It represents packages, classes, and methods in a
Merkle-style hierarchy and replaces application-specific types in method
descriptors with placeholders. This improves resilience to identifier
renaming and permits similarity-based partial matching. Its longitudinal
study also demonstrated the security value of exact library-version
attribution \cite{backes2016libscout}. The main limitation is structural:
package and class hierarchies can change under flattening, shrinking, and
optimization.

LibD moved from names toward internal dependencies and opcode-derived hashes.
Its homogeny graph uses inclusion, inheritance, and call relations to recover
candidate components before feature hashing. The approach was evaluated at
market scale and demonstrated that structural relations can be more robust
than package labels, but it focused on discovering recurring libraries rather
than resolving exact versions \cite{li2017libd}.

ORLIS used a two-stage design: similarity digests first reduced the candidate
space, after which class-level CFG evidence and Jaccard similarity supported a
finer comparison. The work is important because it combined a coarse filter
with partial-library similarity and explicitly studied obfuscation
\cite{wang2018orlis}. Its principal trade-off was computational cost, showing
that structural resilience must be balanced with scalable candidate
reduction.

\subsection{Class Dependencies and Version-Aware Detection}

LibPecker computes strict class signatures from internal class dependencies,
then introduces adaptive class thresholds and weighted similarity to tolerate
customization and eliminated dead code. Its original evaluation reported
98.1\% precision and 91\% recall \cite{zhang2018libpecker}. These results are
methodologically useful, but later common-benchmark evidence shows why
original-paper performance should not be generalized without preserving the
dataset and protocol.

LibID formulates library identification as a constrained correspondence
problem and uses binary integer programming to select consistent matches. It
also uses synthetic applications to tune thresholds and F-Droid applications
as source-backed ground truth. The method identifies specific versions and is
designed to survive code shrinking and package modification
\cite{zhang2019libid}. Its optimization formulation improves global
consistency but increases implementation and runtime complexity.

ATVHUNTER addresses version and vulnerability identification through a
coarse-to-fine pipeline. CFG evidence retrieves potential references, then
opcodes inside basic blocks distinguish candidate versions. The work reports
90.55\% precision and 88.79\% recall and connects detected versions with a
vulnerability database \cite{zhan2021atvhunter}. This is one of the closest
pre-SBOM pipelines to the current project because it links compiled-code
attribution, exact versions, and security consequences. However, its output is
not a standard SBOM, and its evaluation does not fully represent modern R8
method inlining.

\section{Optimization-Aware and Structural Matching}

LibScan builds potential class correspondences and refines them through method
opcode and call-chain opcode similarity. Early-stop rules and reuse of
intermediate results reduce the cost of the deeper comparison. This design
addresses repackaging, dead-code removal, and control-flow changes more
directly than package-only methods \cite{wu2023libscan}.

LibHunter studies a narrower but important question: how compiler
optimizations affect library and version detectors. It explicitly models
inlining and call-site optimization and reports improvements over earlier
tools on optimized, obfuscated, and shrunk applications
\cite{xie2024libhunter}. Its contribution is methodological as well as
algorithmic: optimization must be evaluated separately from identifier
renaming because it can remove or merge the very units used by a matcher.

SAD targets version-level false positives through class structural
correspondence. It first retrieves candidate app classes from class dependency
graphs, compares class summaries, and finally evaluates the structural
similarity of the matched subgraphs. On its obfuscated datasets, SAD reports
97.64\% library-level F1, 84.82\% version-level F1, and 94.12\% class-level F1
\cite{zhou2025sad}. These are strong reported results, but they belong to the
paper's three evaluation datasets and should not be compared directly with a
different real-APK denominator.

The public LibScope artifact represents another direction: replacing
independent local matches with iterative global matching over compact Bloom
filter features \cite{libscopeArtifact}. Global consistency is relevant when
several classes have individually plausible correspondences. At the time of
this review, the accessible artifact was anonymized, so it is included as an
emerging matching direction rather than as a source of quantitative claims.

\section{Learned Representations}

LibAttention replaces hand-written correspondence rules with a feature
language model. APK and library code are converted to Smali, transformed into
an obfuscation-resistant vocabulary, encoded by a Transformer, and compressed
for downstream classification. The paper reports an F1 improvement exceeding
30 percentage points over rule-based methods on its Android R8 dataset
\cite{pan2025libattention}.

The learned approach is valuable because R8 transformations evolve and may not
preserve the exact features expected by a manually designed detector.
Nevertheless, LibAttention performs library-presence classification rather
than exact-version identification. It therefore cannot independently provide
the version evidence required for reliable vulnerability mapping or a
versioned SBOM. Learned representations are better understood here as
assistance for hard identity or risk cases than as automatic version
authority.

\section{SBOM-Oriented Approaches}

Most Android TPL detectors output library names, candidate versions, or class
correspondences rather than a standard SBOM. That difference matters because
an operational SBOM also requires package identifiers, evidence provenance,
component relationships, uncertainty handling, and an interoperable format.

UniBOM is not Android-specific, but it demonstrates a broader SBOM-analysis
architecture that combines multiple evidence sources and supports
\CycloneDX and \SPDX representations \cite{safronov2025unibom}. Its relevance
to this project is architectural: binary recovery should be connected to
standard output, visualization, and vulnerability analysis rather than
stopping at a detector-specific list.

The Android-specific gap remains significant. ATVHUNTER connects versions to
vulnerabilities, and LibScout studies vulnerable library histories, but
neither provides the evidence-qualified \CycloneDX workflow required by this
project. Standard SBOM generation must also avoid converting a top similarity
candidate into an unsupported exact version.

\section{Benchmark and Assessment Evidence}

The 2020 empirical assessment and its extended systematic study compared
publicly available tools under common criteria. They found that no single
method dominated all dimensions: LibScout was strong in effectiveness,
LibRadar in efficiency and usability, and LibPecker in the evaluated
obfuscation settings. More importantly, the studies identified module
decoupling, partial imports, multidex, AAR preprocessing, and tool
complementarity as recurring problems \cite{zhan2020assessment,zhan2021systematic}.

TPL-Benchmark revisits the field with 6,055 applications, 5,756 unique
libraries, 15,274 versions, and 62,218 app-library mappings. Under its unified
protocol, the highest library-level F1 among the ten evaluated tools is
60.15\%, achieved by LibPecker, despite substantially higher values reported
in several original papers \cite{gu2025tplbenchmark}. This does not prove that
the original experiments were incorrect. It demonstrates that performance is
highly sensitive to dataset construction, ground truth, reference coverage,
partial inclusion, and modern compiler transformations.

The benchmark also separates library-level from version-level results and
shows that stronger identity recall does not automatically imply reliable
version selection. This observation directly motivates separate identity and
version confidence in the proposed system.

\section{Method-by-Method Synthesis}

Before comparing the tools in tabular form, this section summarizes the main
idea of each representative approach. The objective is to make clear which
evidence family each method uses and why its result does or does not directly
support an APK-derived SBOM claim.

\textbf{LibScout} builds profiles for known Android libraries and compares an
application against those profiles using package, class, and method
information \cite{backes2016libscout}. Its strength is the explicit reference
model: if the right library version is present in the database and the APK
preserves enough method structure, the match can support a version claim. Its
weakness is that package organization, removed methods, and incomplete
references can reduce the authority of the result.

\textbf{LibD} uses graph-based homogeny relations to identify library code in
large Android application collections \cite{li2017libd}. It is important
because it moves beyond simple package names and uses structural relations
such as inheritance, calls, and inclusion. However, its main role is library
discovery and boundary recovery rather than exact Maven-version reporting.

\textbf{ORLIS} combines coarse and fine matching for obfuscation-resilient
library detection \cite{wang2018orlis}. It first filters candidates using
call-graph-derived features and then verifies class correspondences with fuzzy
signatures. This makes it useful for partial inclusion and obfuscated code,
but the class-mapping result still needs additional evidence before it becomes
a precise SBOM version claim.

\textbf{LibPecker} focuses on class-dependency relations and adaptive
similarity \cite{zhang2018libpecker}. The method is valuable because class
dependencies survive some renaming and package reorganization. Its limitation
is that similarity thresholds and reference coverage still determine whether
the output is a strong identity claim or only a candidate.

\textbf{LibID} formulates library identification as a class-correspondence
problem and uses optimization to find consistent mappings
\cite{zhang2019libid}. This supports version-oriented reasoning better than
package-only methods, especially when code is modified or shrunk. It still
depends on the available reference set and on the transformations represented
in its evaluation.

\textbf{ATVHUNTER} connects version identification to vulnerability analysis
\cite{zhan2021atvhunter}. It decouples candidate library code with class
dependencies, filters with control-flow evidence, and uses opcode-level
matching for version decisions. This makes it directly relevant to security
use cases, but the vulnerability conclusion remains valid only if the
component and version evidence is strong.

\textbf{LibScan} improves fine-grained matching by using class
correspondence, method opcode similarity, and call-chain evidence
\cite{wu2023libscan}. It shows that method-level and call-context evidence
can improve precision, while also showing why a detector must report the
evidence behind each decision instead of exposing only a final label.

\textbf{LibHunter} treats optimization as a first-class problem
\cite{xie2024libhunter}. Its call-site-optimization and cross-inlining
features address transformations that can break simpler method-boundary
assumptions. This is especially relevant for modern Android builds where R8
may inline or rewrite the exact code fragments needed for version
discrimination.

\textbf{SAD} separates stateful and stateless classes and uses different
class-structural evidence for each category \cite{zhou2025sad}. It is one of
the strongest examples of version-level reasoning because it uses field
operations, opcode similarity, class dependencies, and literals. Its
remaining sensitivity to aggressive inlining confirms that version-level
claims still need transformation-aware limits.

\textbf{LibAttention} uses a Transformer model over Smali-oriented
representations for library presence detection \cite{pan2025libattention}.
It demonstrates that learned features can help under obfuscation, but its
output is a presence signal rather than a complete exact-version SBOM entry.
For this project, it is best interpreted as evidence for AI-assisted
candidate support, not as authority to rewrite component versions.

\textbf{UniBOM} is not an Android TPL detector, but it is relevant because it
places binary and source evidence inside an SBOM-oriented workflow
\cite{safronov2025unibom}. It shows the value of combining extraction,
analysis, SBOM generation, and vulnerability visualization. Its limitation
for this thesis is that firmware and source-code SBOM generation do not solve
APK-side Maven exact-version attribution under Android transformations.

\section{Comparative Analysis}

Table~\ref{tab:pfe-sota-comparison} below summarizes the representative tools
using common capability dimensions. Quantitative results remain in the
preceding text so that each number stays attached to its evaluation context.

\begin{landscape}
\thispagestyle{empty}
\noindent
{\small Chapter 4.\quad State of the Art\hfill \thepage\par}
\vspace{1mm}
\hrule
\vspace{5mm}
\begin{table}[H]
\centering
\footnotesize
\caption{Capability-oriented comparison of representative Android TPL and
SBOM approaches.}
\label{tab:pfe-sota-comparison}
\begin{tabularx}{\linewidth}{
  p{0.11\linewidth}
  p{0.17\linewidth}
  p{0.20\linewidth}
  p{0.11\linewidth}
  p{0.15\linewidth}
  X
}
\toprule
\textbf{Approach} & \textbf{Primary goal} & \textbf{Main evidence} &
\textbf{Version level} & \textbf{Optimization focus} & \textbf{Output}\\
\midrule
LibScout & Known library/version matching & Fuzzy method descriptors and
hierarchical hashes & Yes & Renaming; limited structural change &
Detector-specific library/version list\\
LibD & Market-scale library discovery & Dependency relations and opcode
hashing & No & Name obfuscation & Discovered library clusters\\
ORLIS & Obfuscation-resilient matching & CFG structure, sdhash, ssdeep,
Jaccard & Limited & Classical obfuscators & Class/library mappings\\
LibPecker & Precise partial-library matching & Class dependencies, adaptive
thresholds, weighted similarity & Limited & Customization and dead-code
removal & Detector-specific matches\\
LibID & Consistent known-library matching & Class/member features and integer
programming & Yes & Shrinking and package modification &
Library/version identifications\\
ATVHUNTER & Version and vulnerability identification & CFG shortlist and
basic-block opcode evidence & Yes & Obfuscation and dead-code removal &
Version and vulnerability records\\
LibScan & Fine class correspondence & Method and call-chain opcode similarity
& Yes & Repackaging and advanced obfuscation & Detector-specific matches\\
LibHunter & Optimization-aware detection & Inlining and call-site
optimization simulation & Yes & Explicit compiler optimization &
Library/version matches\\
SAD & Version-level structural matching & CDG candidates, class summaries,
matched subgraphs & Yes & Obfuscated applications & Version and class
correspondences\\
LibAttention & Learned library presence & Smali feature language model and
Transformer & No & Android R8 & Binary library presence\\
UniBOM & Unified SBOM analysis & Binary, filesystem, and source evidence &
Depends on input analyzers & Not Android-specific & \CycloneDX and \SPDX\\
\bottomrule
\end{tabularx}
\end{table}
\end{landscape}

The table exposes a recurring trade-off. Fast approaches reduce search cost
but provide less reliable version evidence. Structural and optimization-aware
approaches improve robustness but are more expensive and remain dependent on
reference coverage. Learned methods adapt to transformations but do not
automatically provide explainable version attribution. SBOM-oriented systems
provide interoperable output but need trustworthy component evidence from
upstream analyzers.

\section{Limitations of Existing Approaches}

\subsection{Dataset and Metric Fragmentation}

Original papers frequently evaluate curated libraries or generated
obfuscation variants, while later benchmarks include broader real
applications. Precision, recall, F1, version accuracy, and class
correspondence may also use different units. A value reported for one dataset
cannot be treated as a universal ceiling or directly compared with another
paper's percentage.

\subsection{Package and Reference Dependence}

Package-based segmentation is fragile under renaming and flattening.
Reference-based systems avoid clustering ambiguity but cannot detect a
coordinate or version that is absent from the database. Adding many versions
also increases candidate pressure and can create false near-ties. Reference
coverage and matcher performance must therefore be measured separately.

\subsection{Incomplete R8 and Optimization Coverage}

Renaming is the easiest transformation because much of the bytecode structure
survives. Shrinking removes library code, while optimization and inlining
rewrite method boundaries and call relations. The progression from LibScout
through LibScan, LibHunter, SAD, and LibAttention shows increasing attention
to these transformations, but no reviewed approach establishes universal
exact-version recovery under all modern R8 behaviors.

\subsection{Weak Standard Output and Provenance}

Most detectors stop at tool-specific matches. They do not emit a standard
SBOM, distinguish application files from Maven packages, preserve ambiguous
version sets, or attach the evidence required to audit a claim. Vulnerability
mapping is particularly sensitive to this omission because the same
coordinate at two adjacent versions may produce different advisory results.

\subsection{Limited Uncertainty and Failure Reporting}

A top candidate is not necessarily a confirmed component. Several reviewed
approaches expose similarity scores, but few make identity uncertainty,
version uncertainty, reference absence, and unsupported native or hidden code
first-class output categories. Realistic deployment requires abstention and
bounded claims, not only improved ranking.

\section{Impact on the Proposed MVP}

The project did not copy one detector. It adopted, adapted, or rejected ideas
according to implementation feasibility and target/control evidence.
Table~\ref{tab:pfe-sota-decisions} summarizes the principal decisions.

\begin{table}[H]
\centering
\caption{State-of-the-art ideas and their disposition in the proposed MVP.}
\label{tab:pfe-sota-decisions}
\begin{tabularx}{\textwidth}{L{0.24\textwidth}L{0.15\textwidth}X}
\toprule
\textbf{Research idea} & \textbf{Decision} & \textbf{Project implementation or
rationale}\\
\midrule
Known versioned reference profiles & Adopted & Maven JAR/AAR artifacts are
converted to DEX and stored as coordinate/version feature profiles.\\
Package-only module boundaries & Rejected & Package evidence remains one
signal, but structural, opcode, literal, and dependency evidence are required.\\
Coarse-to-fine candidate verification & Adapted & Candidate reduction and
weighted evidence ranking follow the multi-stage pattern, with fail-fast
controls for unsafe shortlists.\\
Class dependency and call-chain evidence & Adopted diagnostically and
partially in production & Class dependencies contribute to the matcher;
multi-scale features are extractable but remain inert or report-only when
global scoring regresses controls.\\
Generic newest-version tie breaking & Rejected & Recency is not independent
APK evidence and caused target/control conflicts.\\
Learned R8 assistance & Adapted as report-only & Random Forest and Transformer
experiments support risk triage but cannot change selected versions, SBOM
components, or vulnerability labels.\\
Standard SBOM and package identifiers & Adopted & Accepted detections are
serialized as \CycloneDX components with Package URLs and evidence
properties.\\
Single aggregate benchmark score & Rejected & Controlled, clean, partial,
source-known, experiment-only, and limitation lanes are reported separately.\\
\bottomrule
\end{tabularx}
\end{table}

This decision process is as important as feature selection. Opcode n-grams,
API references, multi-scale R8 blends, broad priority databases, and
source-known rerankers were implemented or evaluated, but they remained
disabled, diagnostic, or experiment-only when controls did not justify
promotion.

\section{Research Gap and Positioning}

The reviewed literature leaves five connected gaps:

\begin{enumerate}
  \item exact-version evidence degrades more severely than library identity
  under shrinking and optimization;
  \item reference coverage and matcher recall are often mixed into one score;
  \item standard SBOM generation is rarely integrated with Android binary
  attribution;
  \item vulnerability claims frequently inherit unexamined version
  uncertainty;
  \item negative experiments, abstentions, and boundary cases are not always
  visible in the final result.
\end{enumerate}

The proposed contribution is positioned as an end-to-end,
evidence-gated APK-to-\CycloneDX pipeline. It combines research-backed static
features with versioned Maven references, separates identity from exact
version confidence, validates emitted rows against source when available, and
keeps native, hidden-code, AI, and experimental evidence outside unsupported
production authority.

The objective is not to claim a universally superior detector from
incompatible percentages. It is to provide a working and auditable system in
which each SBOM and vulnerability claim remains connected to its evidence,
denominator, and limitation.

\section{Conclusion}

Android TPL research progressed from fast package and hash signatures to class
dependencies, global correspondence, optimization-aware simulation, and
learned representations. These advances improve robustness, but benchmark
evidence shows that realistic library and version recovery remains difficult,
particularly when reference coverage, partial inclusion, R8 optimization, and
adjacent-version similarity interact.

The literature also reveals a systems gap: accurate detection alone does not
produce a trustworthy SBOM. Standard output, evidence provenance, ambiguity,
source validation, and vulnerability-claim boundaries must be designed
together. Chapter~\ref{ch:pfe-design} answers this gap by presenting the
architecture of the proposed evidence-gated APK-to-SBOM system.
